\begin{table}[ht]
\caption{\textbf{\stackoverflow example}. The positive document provides the documentation that explains the method to solve the problem.}
\centering
% \resizebox{\textwidth}{!}{
\begin{tabular}{p{13cm}}
\toprule
% \hspace{5.5cm} Stackoverflow \\
% \midrule
\textit{\textbf{Query}} \\
\midrule
Is there a Snowflake command that can do the following: \\
\\
a,b,c \\
1,10,0.1 \\
2,11,0.12 \\
3,12,0.13 \\
to a table like this: \\
 \\
key,value \\
a,1 \\
a,2 \\
a,3 \\ 
b,10 \\
b,11 \\
b,13 \\
c,0.1 \\
c,0.12 \\
c,0.13 \\
? \\
 \\
This operation is often called melt in other tabular systems, but the basic idea is to convert the table into a list of key value pairs. \\
 \\
There is an UNPIVOT in SnowSQL, but as I understand it UNPIVOT requires to manually specify every single column. This doesn't seem practical for a large number of columns. \\
\midrule
\textit{\textbf{Chain-of-thought reasoning steps to find documents}} \\
\midrule
This operation is a kind of table transformation like reshaping, serialization, or flattening. We are looking for an operation in Snowflake that can take 2-dimensional values, and transform them into an view that presents one-one correlation mapping between key and value. \\
\midrule
\textit{\textbf{Example positive document}} \\
\midrule
FLATTEN \\
Flattens (explodes) compound values into multiple rows. \\
FLATTEN is a table function that takes a VARIANT, OBJECT, or ARRAY column and produces a lateral view (i.e. an inline view that contains correlation referring to other tables that precede it in the FROM clause). \\
FLATTEN can be used to convert semi-structured data to a relational representation. \\
Syntax \\
FLATTEN( INPUT => <expr> [ , PATH => <constant\_expr> ] \\
\hspace{4.35cm} [ , OUTER => TRUE | FALSE ] \\
\hspace{4.35cm} [ , RECURSIVE => TRUE | FALSE ] \\
\hspace{4.35cm} [ , MODE => 'OBJECT' | 'ARRAY' | 'BOTH' ] )  \\
% Arguments \\
% Required: \\
% INPUT => expr \\
% The expression that will be unseated into rows. The expression must be of data type VARIANT, OBJECT, or ARRAY. \\
% Optional: \\
% PATH => constant\_expr \\
% The path to the element within a VARIANT data structure which needs to be flattened. Can be a zero-length string (i.e. empty path) if the outermost element is to be flattened. \\
% Default: Zero-length string (i.e. empty path) \\
... \\
\midrule
\textit{\textbf{Example negative document}} \\
\midrule
SYSTEM\$EXTERNAL\_TABLE\_PIPE\_STATUS \\
Retrieves a JSON representation of the current refresh status for the internal (hidden) pipe object associated with an external table. \\
Automatically refreshing the metadata for an external table relies internally on Snowpipe, which receives event notifications when changes occur in the monitored cloud storage. For more information, see Introduction to external tables. \\
Syntax \\
SYSTEM\$EXTERNAL\_TABLE\_PIPE\_STATUS( '<external\_table\_name>' ) \\

% MAP\_PICK \\
% Returns a new MAP containing the specified key-value pairs from an existing MAP. \\
% To identify the key-value pairs to include in the new map, pass in the keys as arguments, or pass in an array containing the keys. \\
% If a specified key is not present in the input map, the key is ignored. \\
% Syntax \\
% MAP\_PICK( <map>, <key1> [, <key2>, ... ] ) \\
% MAP\_PICK( <map>, <array> ) \\
% Arguments \\
% <map> \\
% The input map. \\
% <key1>,<key2> \\
% One or more keys that identify the key-value pairs to be included in the returned map. \\
... \\
\bottomrule
\end{tabular}
% }
\label{tab:stackoverflow_example}
\end{table}